\section{Localized modes and self-adjoint eigenproblems}
\label{sec:localized}

A core goal of the substrate theory is to explain particle-like behavior (localization, discrete spectra,
internal ``spin/charge''-like structure) as properties of confined wave states rather than point particles.
In this framework, particles are extended standing-wave configurations of the substrate.

\subsection{Linearization about symmetric backgrounds and self-adjointness}

Let $\vecX_0(x)$ be a time-independent background embedding representing a localized deformation region.
We are especially interested in backgrounds that are \emph{spherically symmetric} in the intrinsic brane coordinates,
so that invariants depend only on $r=|x|$.

Linearize the equations of motion \eqref{eq:eom} about $\vecX_0$ by writing
$\vecX(x,t)=\vecX_0(x)+\bm{\eta}(x,t)$ with small perturbation $\bm{\eta}$.
Because the dynamics derives from the action \eqref{eq:action}, the linearized operator is
the second variation of $S$ and is therefore self-adjoint with respect to the $\rho_m$-weighted
$L^2$ inner product on $\Omega$; orthogonality and completeness of eigenmodes follow as standard consequences.

\subsection{Angular structure: scalar and vector spherical harmonics}
\label{sec:angular-structure}

The substrate is the 6-neighbor cubic lattice, whose point group is $O_h$ rather than the full $SO(3)$. The
long-wavelength continuum recovers approximate isotropy at scales $\gg a$, with residual cubic anisotropy of order
$7.5\%$ in the longitudinal speed at the operating point $\alpha = 0.2$. Solitons in this theory live
in the Skyrme-stabilized regime of width $\gg a$, so the emergent approximate $SO(3)$ is the right symmetry to
organize their angular structure, with $O_h$ corrections entering as small splittings of $SO(3)$ multiplets.

For a spherically symmetric background, the linearized operator commutes with the action of $SO(3)$ on the
intrinsic coordinates at leading order in $a/\ell_{\mathrm{soliton}}$.
Consequently, eigenmodes can be chosen to transform in irreducible $SO(3)$ representations.
For a scalar field this gives the familiar separation into a radial problem and an angular problem whose
eigenfunctions are the spherical harmonics $Y_{\ell m}(\theta,\varphi)$.

The lateral triplet $\xi^i \in \mathbb{R}^3$, and its complex-envelope promotion $\Psi \in \mathbb{C}^3$ from
\Cref{sec:geophase}, carry an internal vector index $i$ in addition to spatial dependence. The correct angular
basis is therefore the \emph{vector spherical harmonic} basis $\mathbf{Y}^{i}_{JLM}$, where $J$ is the total
angular momentum of the field, $L$ is the orbital angular momentum, and $L \in \{J-1,\, J,\, J+1\}$:
\begin{equation}
  \xi^{i}(r,\theta,\varphi) \;=\; \sum_{J,L,M} f_{JLM}(r)\, \mathbf{Y}^{i}_{JLM}(\theta,\varphi).
  \label{eq:vsh-decomposition}
\end{equation}
The choice of $(J,L)$ specifies how the internal index $i$ is glued to the spatial angular coordinates --- it is
the choice of how the lateral triplet's $U(3)$ gauge structure (\Cref{sec:geophase}) couples to the geometry of
the localized mode. Configurations with $J = 0$ are spherically symmetric only under the \emph{combined} rotation
that rotates both spatial and internal indices together; this combined rotation is the physically meaningful
notion of ``spherical symmetry'' for a triplet field.

At soliton scale, $O_h$ refinements enter as small splittings of $SO(3)$ multiplets.
The $L = 1$ representation maps intact to the $T_{1u}$ irrep of $O_h$, so vector ansätze of total orbital
$L = 1$ survive without restriction.
Higher-$L$ modes split (e.g.\ $L = 2 \to E_g \oplus T_{2g}$) and should be expressed in cubic harmonics whenever
lattice corrections matter.

This framework is descriptive, not structural: it organizes the angular content of emergent localized modes, but
it does not determine which modes are physically realized. Which $(J,L)$ channels support stable stationary
solutions is decided by the substrate dynamics together with the geometric quartic stabilization mechanism
discussed in \Cref{sec:continuum-model} --- whose lattice-exact coefficient is $\propto k_s\alpha/a$
(vanishing as $\alpha\to 0$), so soliton stabilization requires $\alpha$ significantly above zero.

Appropriate radial superpositions, together with nontrivial embedding-space polarization structure, can yield
torus-like, loop-like, or shell-like energy density profiles.
These should be understood as standing-wave organizations of the substrate rather than as literal point particles
moving on thin trajectories.

\subsection{Combined particle labels at the soliton layer}
\label{sec:soliton-labels}

A localized 3D mode is fully classified at the soliton layer by the labels
\begin{equation}
  \big( J,\,P;\;\; \mathrm{SU}(3)\text{-irrep};\;\; Q_{\mathrm{U}(1)};\;\; B_{\mathrm{winding}} \big),
  \label{eq:soliton-labels}
\end{equation}
where:
\begin{itemize}
  \item $J$ is the total angular momentum from the emergent approximate $SO(3)$ at scales $\gg a$. On the cubic
    lattice this label refines to an $O_h$ irrep at sub-leading order in $a/\ell_{\mathrm{soliton}}$.
  \item $P$ is parity, the discrete element of the emergent rotation group.
  \item The $\mathrm{SU}(3)$ irrep and $\mathrm{U}(1)$ charge come from the $U(3)$ decomposition of the lateral
    triplet (\Cref{sec:geophase}).
  \item $B_{\mathrm{winding}}$ is the topological winding number of the field configuration.
\end{itemize}

These labels are not independent: the choice of how the internal triplet index couples to the spatial angular
variable via the vector spherical harmonic basis sets a Clebsch--Gordan relationship between the $SO(3)$ and
$\mathrm{U}(3)$ representations carried by the field.
The canonical baryon ansatz is the \emph{hedgehog}
\begin{equation}
  \xi^{i}(\vec{x}) \;=\; f(r)\, \hat{x}^{i},
  \label{eq:hedgehog}
\end{equation}
with $J = 0$, $L = 1$.
Here the internal index $i$ is locked to the spatial direction $\hat{x}^i$, so the configuration is invariant only
under the combined diagonal $SO(3)$.
The entire nontrivial structure of the hedgehog lives in the $\mathrm{SU}(3)$ traceless sector: the $\mathrm{U}(1)$
trace $f(r)\,(\hat{x}^1 + \hat{x}^2 + \hat{x}^3)$ averages to zero on the sphere, giving a trace-neutral far field.
Topological stabilization (winding $B = 1$) is achieved by extending the hedgehog into a Skyrme-twisted
configuration in which the embedding-amplitude direction $X^4$ also rotates radially as a function of $r$.

The proton-versus-neutron distinction enters as a $J = 0$ trace admixture on top of the hedgehog: a nonzero
scalar $L = 0$ component of $\xi^i$ shifts the $\mathrm{U}(1)$ far field to proton-like, while a pure hedgehog or
Skyrme-twisted hedgehog has neutron-like far field.

\subsection{Amplitude scaling and energy normalization}

The linear eigenproblem fixes mode shapes and frequencies but not the overall amplitude:
if $\bm{\eta}$ is an eigenmode, so is $c\,\bm{\eta}$ for any scalar $c$.
Physical amplitudes are selected only when nonlinearities are included and the total energy is fixed.
For an electron-like excitation one may impose a rest-energy normalization
\begin{equation}
  \overline{E}[\vecX] \;=\; m_e c^2,
  \label{eq:rest-energy-normalization}
\end{equation}
where $\overline{E}$ denotes the conserved substrate energy associated with \eqref{eq:action}.
For baryon-like modes the same logic applies with the appropriate target rest energy of the chosen species.
The important structural point is unchanged: amplitude is not determined by the linear eigenproblem alone.

\subsection{Confinement as geometric self-guidance}

A confining mechanism is required: in a purely linear wave system, localized packets generally disperse or radiate.
A useful conceptual benchmark is the toroidal electron proposal attributed to Williamson \& van der Mark, which argues
that ordinary gravitational curvature is far too weak to confine a Compton-scale photon-like mode and that confinement
must arise from a self-consistent, excitation-specific mechanism.
We adopt the \emph{confinement requirement} while not committing to a literal double-loop trajectory
(see e.g.\ the review discussion in \citep{HuygensOptics2025}).

In the present substrate theory, the first candidate confinement mechanism is the brane's \emph{geometric}
self-guidance.
Large-amplitude deformations modify the local Euclidean link lengths on the lattice and the induced metric in the
continuum; this changes the local stress response and therefore the propagation characteristics of fluctuations around
the deformed state.
The crucial point is that this mechanism is already present even when the microscopic link law is Hooke-linear in the
scalar extension.
Constitutive refinements may strengthen or weaken localization, but the underlying nonlinear feedback comes from the
same geometry that couples transverse amplitude to lateral strain.

\subsection{Kinematic confinement of color}
\label{sec:color-confinement}

Before asking whether the substrate supports a bound colored mode, it is worth recording a
structural fact about the \emph{unbound} spectrum: the linear theory carries no colored asymptotic
free state. Recall from \Cref{sec:geophase} that the $U(1)$ trace of the lateral triplet is read as
electromagnetism and the $SU(3)$ traceless content as color, both carried by the per-wavepacket
complex envelope $\Psi \in \mathbb{C}^3$.

On the 6-neighbor stencil the dynamical matrix $D(k)$ is diagonal in the axis-aligned basis, with
three lateral branch frequencies $h_a(k)$ whose spread is measured by the directional anisotropy
$g(\hat k) = \sqrt{\sum_a (h_a - \bar h)^2}/\bar h$. Two requirements on a freely propagating triplet
wavepacket are then mutually exclusive:
\begin{itemize}
  \item \textbf{Coherence.} A wavepacket whose three components stay phase-aligned as it propagates
    requires the three branches to be degenerate, $g(\hat k) = 0$. On the cubic lattice this holds
    \emph{only} along the body-diagonal $k \parallel [111]$.
  \item \textbf{Color-activity.} Operationally meaningful $SU(3)$-traceless content --- a finite
    fibre-internal holonomy distinguishing the three axis components --- requires the branches to be
    \emph{non-degenerate}, $g(\hat k) \neq 0$, i.e.\ propagation \emph{off} $[111]$. Exactly on
    $[111]$ the degenerate triplet has no preferred internal basis and the color holonomy is
    undefined.
\end{itemize}
No linear propagation direction satisfies both at once: off $[111]$ the unequal branch speeds
dephase the triplet over long range, while on $[111]$ there is no resolvable color. Consequently
\emph{color charge cannot be carried to infinity by a free linear excitation.} It has support only
in a nonlinear, phase-locked bound state, where the three components are held together by the
geometric coupling of \Cref{sec:continuum-model} rather than by spectral degeneracy. We refer to
this as the \emph{kinematic confinement of color}.

This is a statement about asymptotic states --- there is no colored free particle --- and not a
dynamical derivation of a linear inter-quark potential, string tension, or area law, which remain
outside the present scope (\Cref{sec:scope}). It is consistent with the vanishing of the
Brillouin-zone Berry/Wilczek--Zee curvature noted in \Cref{sec:geophase}: color is not a curvature
on the Brillouin zone but a fibre-internal object on $(x,t)$ that only closes around a localized
worldtube. The picture carries a falsifiable, normalization-independent handle: the off-$[111]$
fibre-holonomy ratio scales with prestress as
$R(\alpha_2)/R(\alpha_1) = \big[(3-2\alpha_2)/\alpha_2\big]\big/\big[(3-2\alpha_1)/\alpha_1\big]$,
equal to $0.3077$ for $(\alpha_1,\alpha_2) = (0.2,0.5)$; a deviation above $\sim 10\%$ would falsify
the spectral-susceptibility factorization underlying it.

\subsection{Two complementary numerical targets}
\label{subsec:baryon-next-target}

The theory contains two distinct particle sectors with different topology and therefore different numerical
questions; neither has yet been demonstrated as a stable self-consistent solution.

\paragraph{Target 1: the $U(1)$ carrier-phase vortex (electromagnetic / electron-like sector).}
The $U(1)$ trace sector supports line vortices ($\pi_1(U(1))=\mathbb{Z}$), appearing as codimension-2 objects in 3D.
However, once the $SU(3)$ sector is unfrozen the vacuum manifold is simply connected, so this vortex is
\emph{not} topologically protected: it is a semilocal vortex \citep{VachaspatiAchucarro1991Semilocal,AchucarroVachaspati2000}, stable only as a dynamical bound state rather than by topology.
A preliminary rotating-frame-periodic JFNK relaxation of a $Y_1^1$ pure-trace vortex seed on a $112^3$ lattice did not converge: the core drifted $\approx 19$ lattice units off the symmetry axis, the ring amplitude became azimuthally asymmetric, and the measured winding changed sign relative to the seed (\Cref{fig:seed-snapshot}).
This is the expected signature of semilocal orientational moduli \citep{Eto2021OrientationalModuli}: with a dynamical $SU(3)$ sector, the $U(1)$ vortex is free to drift.
A converged eigenstate has not been obtained; branch-selection machinery for stabilizing the semilocal vortex remains future work.

\begin{figure}[tbp]
  \centering
  \includegraphics[width=\textwidth]{diagrams/exp_relaxed_snapshot}
  \caption{Preliminary $U(1)$-trace vortex candidate after 12 rotating-frame-periodic JFNK relaxation steps on a $112^3$ lattice (residual per DOF $\approx 1.9\times10^{-4}$; not a converged eigenstate).
  \textit{Left:} XZ midplane, carrier phase as RGB.
  \textit{Centre:} XY midplane plan view.
  \textit{Right:} XZ midplane amplitude, showing core drift and azimuthal asymmetry characteristic of semilocal-vortex orientational moduli.
  No claim is made of a stable soliton or emergent electromagnetism.}
  \label{fig:seed-snapshot}
\end{figure}

\paragraph{Target 2: the baryon / $\pi_3$ texture sector.}
The cubic substrate naturally singles out three axis-associated channels, and $\pi_3(SU(3))=\mathbb{Z}$ admits
topologically stable textures --- Skyrmion-like lumps --- distinct from the vortex sector.
That makes the proton/neutron sector a particularly informative target: a confined three-component standing-wave
mode tests the triplet coupling structure that the lattice introduces most directly, and topological winding
provides a non-trivial stabilization mechanism not available to the $U(1)$ vortex.
In contrast, an isolated electron-like ansatz --- lacking topological protection --- gives less direct leverage
on the triplet mixing sector by itself.

The ansatz labels (hedgehog \eqref{eq:hedgehog}, Skyrme-twisted extension, trace admixture for proton/neutron
distinction) are defined in \Cref{sec:soliton-labels}.
The concrete four-candidate ansatz menu used in the simulation search is built from these labels.
These are concrete questions for simulation, not derivations of QCD or nucleon phenomenology.
For each target, the first binary question is: does the substrate support a stable, localized solution at all?
If the answer is no for either, an important bridge in the theory fails early and clearly;
if yes, the respective sector becomes the natural arena for sharpening the particle interpretation.

\subsection{Connection to Berry/Wilczek--Zee holonomy and spin}

Localized modes can carry internal polarization structure (multi-component standing waves).
The Berry/Wilczek--Zee connections of \Cref{sec:geophase} can then be evaluated on loops surrounding the localized
region or on adiabatic deformations of a localized mode family.
In this way, long-range ``charge''-like labels are interpreted as far-field holonomy/phase structure, while the
spin-$\tfrac12$ $4\pi$ aspect is attributed to spinorial holonomy of a two-level polarization subspace.
